% Copyright 2026 Benjamin Stanley Frohman (ORCID 0009-0006-7068-3718). Apache License 2.0.
\documentclass[11pt]{article}
\usepackage{amsmath,amssymb,amsthm}\usepackage[utf8]{inputenc}
\newtheorem{theorem}{Theorem}\newtheorem{lemma}[theorem]{Lemma}
\theoremstyle{remark}\newtheorem*{commentary}{Commentary}
\title{Fourier Lemmas on the Three-Torus: Heat Smoothing, the $L^2$ Strain Identity,\\ and the Curl of the Energy Derivative, Formalised in Lean~4}
\author{Benjamin Stanley Frohman\\ORCID 0009-0006-7068-3718}
\date{DRAFT --- 30 October 2026}
\begin{document}\maketitle
% The mathematics in every section is taken from the Lean build (TorusFourierNS/*.lean).
% Commentary blocks quote the author's earlier drafts and notes, with the source and date.

\begin{abstract}
On $\mathbb T^3=(\mathbb R/\mathbb Z)^3$, the product of three circles, we prove in Lean~4 with Mathlib,
with no \texttt{sorry} and only the standard axioms: the heat semigroup is the Fourier multiplier
$e^{-\nu|n|^2t}$, obeys the semigroup law and contracts $L^2$; it maps $H^s$ to $H^{s+k}$ with
$\|e^{t\nu\Delta}f\|_{H^{s+k}}^2\le k!\,e^{2\nu t}(2\nu t)^{-k}\|f\|_{H^s}^2$; for divergence-free
velocity fields $\|S\|_2^2=\tfrac12\|\omega\|_2^2$; and the kinetic energy, viewed as a function of
vorticity, has derivative the stream function, whose curl is the velocity.
\end{abstract}

\section{Setting}
$\mathbb T^3$ carries normalised Haar measure; $e_n(x)=e^{2\pi i n\cdot x}$, $n\in\mathbb Z^3$, is a
Hilbert basis of $L^2(\mathbb T^3)$ (Mathlib's \texttt{mFourierBasis}), and
$\hat f(n)=\int f\,\overline{e_n}$.
\begin{lemma}[Parseval, \texttt{parseval\_norm\_sq}] $\|f\|_2^2=\sum_n|\hat f(n)|^2$.\end{lemma}
\begin{proof} The coefficient map is a linear isometry $L^2\to\ell^2(\mathbb Z^3)$ and
$\|a\|_{\ell^2}^2=\sum|a_n|^2$.\end{proof}
\begin{lemma}[\texttt{norm\_multiplier\_le}] If $|m(n)|\le1$ for all $n$, then $T_m f=\sum m(n)\hat f(n)e_n$
defines $T_m:L^2\to L^2$ with $\|T_mf\|_2\le\|f\|_2$.\end{lemma}
\begin{proof} $\sum|m(n)\hat f(n)|^2\le\sum|\hat f(n)|^2<\infty$, so the sequence lies in $\ell^2$;
pull back by the isometry and apply Parseval.\end{proof}
\begin{commentary} The earlier manuscript already ran the analysis through Fourier coefficients:
``The Fourier multiplier proof of the Sobolev embedding $H^s(\mathbb T^3,\mathbb R^3)\hookrightarrow
C^1(\mathbb T^3,\mathbb R^3)$ for $s>5/2$ follows from Parseval's identity and the $p$-series
convergence'' (\emph{vFinal}, Drive). When this build was planned, the author's point was
simply: ``the torus is proved to be a Circle (cubic)'' (30 October 2026). $\mathbb T^3$ is three
circles, so Mathlib's multi-dimensional Fourier series and Parseval apply directly, and the Sobolev
scale, missing from Mathlib, is defined here as a weighted coefficient sum.\end{commentary}

\section{Heat semigroup}
\begin{theorem} For $\nu,s,t\ge0$: $\widehat{e^{t\nu\Delta}f}(n)=e^{-\nu|n|^2t}\hat f(n)$,
$e^{0}=\mathrm{id}$, $e^{s\nu\Delta}e^{t\nu\Delta}=e^{(s+t)\nu\Delta}$, $\|e^{t\nu\Delta}f\|_2\le\|f\|_2$.\end{theorem}
\begin{proof} The symbol has modulus $\le1$; $e^{-\nu|n|^2s}e^{-\nu|n|^2t}=e^{-\nu|n|^2(s+t)}$ and
operators with equal coefficients are equal.\end{proof}

\begin{commentary} The local-existence block of the May drafts used exactly this operator:
``The heat semigroup $e^{\nu t\Delta}$ is bounded on $H^s$ ($s>5/2$) with
$\|e^{\nu t\Delta}\|_{\mathcal L(H^s)}\le1$'' (\emph{May 16 edits}, Drive). There it was asserted;
here it is the contraction theorem above, for every $s$, since the symbol has modulus at most one
at every mode.\end{commentary}

\section{Sobolev smoothing}
$\|f\|_{H^s}^2=\sum_n(1+|n|^2)^s|\hat f(n)|^2$.
\begin{lemma} For $k\in\mathbb N$, $a>0$, $x\ge0$: $(1+x)^ke^{-ax}\le k!\,e^a a^{-k}$.\end{lemma}
\begin{proof} $y^k/k!\le e^y$ with $y=a(1+x)$, and $e^{-ax}=e^ae^{-a(1+x)}$.\end{proof}
\begin{theorem}[\texttt{heat\_smoothing}] For $\nu,t>0$, $s\in\mathbb R$, $k\in\mathbb N$, $f\in H^s$:
$e^{t\nu\Delta}f\in H^{s+k}$ and $\|e^{t\nu\Delta}f\|_{H^{s+k}}^2\le k!\,e^{2\nu t}(2\nu t)^{-k}\|f\|_{H^s}^2$.\end{theorem}
\begin{proof} Mode by mode,
$(1+|n|^2)^{s+k}e^{-2\nu t|n|^2}|\hat f|^2\le k!e^{2\nu t}(2\nu t)^{-k}(1+|n|^2)^s|\hat f|^2$ by the lemma with
$a=2\nu t$; sum by comparison.\end{proof}
\begin{theorem}[\texttt{heat\_smoothing\_real}] The same holds for every real $k\ge0$ with $k!$
replaced by $\lceil k\rceil!$.\end{theorem}
\begin{proof} For $y\ge0$, $y^k\le\lceil k\rceil!\,e^y$: if $y\le1$ then $y^k\le1\le e^y$; if $y\ge1$
then $y^k\le y^{\lceil k\rceil}\le\lceil k\rceil!\,e^y$. The rest is as before.\end{proof}
\begin{commentary} The v2.4 master draft names this step as indispensable: ``Viscosity is essential
for local existence in $H^s$, the explicit higher-norm Gronwall estimate, and the final parabolic
smoothing'' (\emph{NS Global Regularity v2.4 Master}, Drive, 25 May 2026). The drafts invoked the
smoothing without a statement; the theorem above gives it with an explicit constant.\end{commentary}

\section{Strain and vorticity}
For $k\in\mathbb R^3$, $\hat u\in\mathbb C^3$, $k\cdot\hat u=0$: $\hat S_{ij}=\pi i(k_i\hat u_j+k_j\hat u_i)$, $\hat\omega=2\pi i\,k\times\hat u$.
\begin{lemma} $\sum_{ij}|\hat S_{ij}|^2=\tfrac12|\hat\omega|^2$ and $|\hat\omega|^2=4\pi^2|k|^2|\hat u|^2$.\end{lemma}
\begin{proof} Splitting $\hat u=a+ib$, for each real part
$\sum_{ij}(k_ia_j+k_ja_i)^2=2(|k|^2|a|^2+(k\cdot a)^2)$ and $|k\times a|^2=|k|^2|a|^2-(k\cdot a)^2$; both
$k\cdot a$ and $k\cdot b$ vanish.\end{proof}
\begin{theorem}[\texttt{L2\_strain\_eq}] For $L^2$ fields with these coefficients, $\|S\|_2^2=\tfrac12\|\omega\|_2^2\le\|\omega\|_2^2$.\end{theorem}
\begin{proof} Parseval on each component, exchange the finite and infinite sums, apply the lemma per mode.\end{proof}

\section{The curl of the energy derivative}
$\hat\psi=\hat\omega/(4\pi^2|k|^2)$, $H(\hat\omega)=|\hat\omega|^2/(8\pi^2|k|^2)$.
\begin{theorem} For $k\ne0$, $k\cdot\hat u=0$: $H(\mathrm{curl}\,\hat u)=\tfrac12|\hat u|^2$;
$H(\omega+\varepsilon\eta)=H(\omega)+\varepsilon\,\mathrm{Re}\langle\hat\psi,\eta\rangle+\varepsilon^2H(\eta)$, so
$\delta H/\delta\omega=\hat\psi$; and $2\pi i\,k\times\hat\psi=\hat u$, i.e.\ $\mathrm{curl}(\delta H/\delta\omega)=u$.\end{theorem}
\begin{proof} $k\times(k\times\hat u)=k(k\cdot\hat u)-|k|^2\hat u=-|k|^2\hat u$, so
$(2\pi i)^2k\times(k\times\hat u)/(4\pi^2|k|^2)=\hat u$. The expansion is the polarisation of a real quadratic form.\end{proof}
\begin{commentary} The original degeneracy argument read: ``2. $\delta H/\delta\omega = u =$
Biot--Savart $\omega$. 3. By definition of $\Pi_u$ we have $\Pi_u(u)=0$'' (\emph{Tether\_May31.lean},
Drive), and the May drafts wrote ``$\delta H/\delta\mu=u$ (Biot--Savart recovery, Coulomb gauge)''
(\emph{May 16 edits}). Formalising it exposed the slip: as a functional of vorticity, the derivative
of the kinetic energy is the stream function $\psi$, not $u$, and the $\psi$-slot tether failed energy
degeneracy for that reason. The repair was to put the curl into the tether slot, so that what it
needs is $\mathrm{curl}(\delta H/\delta\omega)=u$. In the tether build that identity was a hypothesis,
because the derivative there is abstract. In Fourier variables it is the theorem above: the earlier
statement becomes true once the curl is put back in.\end{commentary}

\section{Scope}
Only $p=2$ strain bounds are proved. Nothing here concerns
existence or regularity of Navier--Stokes solutions.
\end{document}
